\section{引言：数据分析的核心目标}

数据分析的本质是\textbf{用数据驱动决策}。目标不是分析本身，而是产出一个能帮助他人行动的 artifact（交付物）：为管理层提供的图表、为产品团队提供的实验报告、为研究者提供的模型评估、或指导日常运营的 dashboard。

\begin{knowledgebox}{数据分析循环框架}
《R for Data Science》一书推广了一种经典的数据分析循环：
\begin{enumerate}[nosep]
    \item \textbf{Import}（导入）：加载原始数据
    \item \textbf{Tidy}（整理）：清洗、规范化数据格式
    \item \textbf{Transform / Visualize / Model}（变换 / 可视化 / 建模）：在三者之间迭代，逐步建立对数据的理解
    \item \textbf{Communicate}（沟通）：将结果包装为他人可消费的形式
\end{enumerate}
编程贯穿整个循环。Codex 的价值在于加速你在循环中的每一步。
\end{knowledgebox}

\begin{importantbox}{Codex 的定位}
Codex 不是用来生成一次性 notebook 的工具。它的目标是帮助你构建一个\textbf{其他人可以审查、信任并重新运行}的完整工作流（reproducible workflow）。
\end{importantbox}

\subsection{本章小结}
数据分析应当以``产出可行动的交付物''为导向，Codex 帮助你在 import--tidy--transform--visualize--model--communicate 循环中更高效地推进。

